\documentclass[lettersize,journal]{IEEEtran}

\usepackage{amsmath,amsfonts,amssymb}
\usepackage{algorithmic}
\usepackage{algorithm}
\usepackage{array}
\usepackage[caption=false,font=normalsize,labelfont=sf,textfont=sf]{subfig}
\usepackage{textcomp}
\usepackage{stfloats}
\usepackage{url}
\usepackage{verbatim}
\usepackage{graphicx}
\usepackage{booktabs}
\usepackage{cite}

\hyphenation{op-tical net-works semi-conduc-tor IEEE-Xplore}

\begin{document}

\title{Deadline and Reliability Constrained Hierarchical Reinforcement Learning
for Energy-Efficient Dynamic Edge--Cloud Workflow Scheduling}

\author{%
% Author information will be added before submission.
}

% Add the target journal header after the venue is selected.
% \markboth{Journal Name, Vol. XX, No. XX, Month Year}%
% {Author \MakeLowercase{\textit{et al.}}: Short Paper Title}

\maketitle

\begin{abstract}
% To be completed after the method and experimental results are finalized.
\end{abstract}

\begin{IEEEkeywords}
% To be completed.
\end{IEEEkeywords}

\section{Introduction}
\label{sec:introduction}

\IEEEPARstart{E}{merging} applications such as industrial Internet of Things, intelligent
healthcare, and distributed data analytics increasingly execute close to data
sources while retaining access to elastic cloud resources. Edge--cloud
collaboration is therefore attractive because edge servers shorten data paths
and cloud servers provide greater computing capacity. Many such applications
comprise multiple dependent stages rather than independent requests. They are
naturally represented as directed acyclic graph (DAG) workflows, in which
vertices are computational tasks and edges encode precedence relations and
inter-task data transfers~\cite{topcuoglu2002,jayanetti2022}. Scheduling these
workflows is difficult in a heterogeneous edge--cloud system because a task's
completion time depends jointly on predecessor completion, communication
delay, the selected server, and the server's current queue.

The difficulty is amplified in an online multi-workflow setting. Workflows
arrive over time, share the same servers, and compete for queue positions.
Consequently, placing one ready task changes not only its own execution time
but also the congestion observed by tasks that arrive or become ready later.
Classical list scheduling, exemplified by HEFT~\cite{topcuoglu2002}, provides
an effective static ordering and assignment principle when the workflow and
resource state are known in advance. Dynamic workflow schedulers and
evolutionary methods extend this setting to changing arrivals, objectives, or
resources~\cite{arabnejad2019,ismayilov2020}, but repeated search or
replanning becomes increasingly costly as the numbers of concurrent workflows
and heterogeneous servers grow.

Energy management introduces two coupled service-quality risks. Dynamic
voltage and frequency scaling (DVFS) can reduce processor dynamic power, but a
lower frequency also prolongs execution, consumes the remaining critical-path
slack, and propagates delay through shared queues. Moreover, voltage and
frequency reduction can increase the transient-fault rate, so the most
energy-efficient decision is not necessarily reliable. Previous
reliability-aware workflow schedulers have addressed this trade-off through
server remapping, task replication, checkpointing, or reliability-aware
frequency selection~\cite{zhang2017,wu2018,hassan2020,wu2023}. These mechanisms
show that energy, deadline satisfaction, and reliability must be considered
jointly rather than optimized in isolation.

Deep reinforcement learning (DRL) provides a promising way to adapt scheduling
decisions to online system states. Existing studies learn task placement for
dynamic or heterogeneous workflows and optimize makespan, utilization, cost,
or energy through fixed reward terms or learned objective weights
~\cite{jayanetti2022,chen2023,pan2024,wen2026}. Nevertheless, three limitations
remain for the problem considered here. First, a server choice is discrete,
whereas its DVFS ratio is continuous; flattening both decisions into one action
space obscures their different physical roles and enlarges the search space.
Second, deadline and reliability requirements are commonly represented by
fixed penalties or best-effort objectives, whose effective trade-offs can
change with workload and system scale. Although constrained reinforcement
learning offers adaptive Lagrangian mechanisms
~\cite{achiam2017,tessler2019,stooke2020}, directly assigning workflow-level
violations to individual
scheduling decisions remains difficult. Third, a local DVFS decision can delay
other workflows through a server queue, so a terminal miss signal alone does
not identify whether the dominant cause was server placement or frequency
reduction.

To address these challenges, this paper studies computing-energy minimization
for dynamically arriving DAG workflows under workflow-level deadline and
reliability requirements. We develop a constraint-aware hierarchical
reinforcement learning framework. The upper level uses a three-head Dueling
Double DQN~\cite{vanhasselt2016,wang2016dueling} to select a server. Its heads
estimate energy reward, deadline cost, and reliability cost separately, and
online Lagrange multipliers combine the three values only when an action is
selected. Conditioned on the chosen server, the lower level uses Soft
Actor--Critic (SAC)~\cite{haarnoja2018} to produce a continuous DVFS ratio.
This decomposition matches the mixed action structure while allowing the two
levels to use different state information and constraint feedback.

The framework further introduces three mechanisms for long-horizon constraint
credit. Queue-congestion credit attributes the waiting time induced by a
placement decision to the upper policy. A remaining-critical-path safety state
describes how much execution time can be assigned to the current task without
exhausting the estimated downstream slack. At the lower level, a workflow-level
counterfactual replay holds task placements and within-server orders fixed,
restores the source workflow to maximum frequency, and propagates precedence
and queue constraints again. The resulting change in deadline cost estimates
the responsibility of DVFS, including its effects on other workflows sharing
the same queues. The upper and lower levels maintain separate deadline
multipliers, while the lower level additionally maintains a reliability
multiplier for frequency-dependent failures.

The main contributions of this paper are summarized as follows:
\begin{enumerate}
    \item We formulate online scheduling in a heterogeneous edge--cloud system
    as a mixed discrete--continuous optimization problem that minimizes
    computing energy under workflow-level deadline and reliability
    requirements, stochastic workflow arrivals, precedence constraints, and
    shared non-preemptive queues.

    \item We design a hierarchical constrained policy in which a three-head
    Dueling Double DQN performs discrete server placement and a continuous-action
    SAC controller determines the DVFS ratio. Independent deadline multipliers
    at the two levels and a lower-level reliability multiplier adapt constraint
    pressure without relying on fixed multi-objective weights.

    \item We develop queue-congestion credit, a remaining-critical-path safety
    representation, and workflow-level DVFS counterfactual attribution to
    connect delayed workflow violations to the placement and frequency
    decisions that caused them.

    \item We adopt constraint-first checkpoint selection based on the empirical
    success probability at a 90\% deadline-success threshold (ESP@90) and the
    reliability success rate. Energy is compared only among checkpoints that
    pass the prescribed validation gates.
\end{enumerate}

The proposed approach targets empirical constraint satisfaction over the
workflow distribution. It does not claim a deterministic hard-real-time or
per-workflow reliability guarantee for every unseen arrival sequence.

\section{Related Work}
\label{sec:related_work}

\subsection{Dynamic Workflow Scheduling in Heterogeneous Edge--Cloud Systems}

\subsection{Constrained and Hierarchical Reinforcement Learning}

\subsection{DVFS and Reliability-Aware Scheduling}

\section{Problem Formulation}
\label{sec:problem_formulation}

We consider an online edge--cloud platform that receives a stream of DAG
workflows. To keep the formulation concise, this section describes an
arbitrary arriving workflow $G_k$; the same model applies to every workflow in
the stream. Contention with earlier and concurrent workflows is captured by
the availability time of each server. The main notation is summarized in
Table~\ref{tab:key_notations}.

\begin{table}[t]
    \centering
    \caption{Key Notations}
    \label{tab:key_notations}
    \renewcommand{\arraystretch}{1.08}
    \begin{tabular}{@{}p{0.25\columnwidth}p{0.68\columnwidth}@{}}
        \toprule
        Notation & Description \\
        \midrule
        $G_k$ & workflow $k$, described by $(\mathcal{N}_k,\mathcal{E}_k,D_k,\Theta_k)$ \\
        $\mathcal{M}$ & set of edge and cloud servers \\
        $w_i,d_{ij}$ & workload of task $i$ in Giga-cycles and data on edge $(i,j)$ \\
        $f_m^{\max}$ & maximum processing frequency of server $m$ in GHz \\
        $m_i,\rho_i$ & selected server and DVFS ratio of task $i$ \\
        $RT_i,AST_i,AFT_i$ & ready, actual-start, and actual-finish times \\
        $C_k,E_k,R_k$ & response time, computing energy, and reliability of workflow $k$ \\
        \bottomrule
    \end{tabular}
\end{table}

\subsection{System and Scheduling Model}
\label{sec:system_model}

\subsubsection{Workflow and Edge--Cloud Resource Model}

Each workflow $G_k$ is represented by a DAG
\begin{equation}
    G_k=(\mathcal{N}_k,\mathcal{E}_k,D_k,\Theta_k),
    \label{eq:workflow_model}
\end{equation}
where $\mathcal{N}_k$ is the task set, $\mathcal{E}_k$ is the set of
precedence edges, $D_k$ is the deadline, and $\Theta_k$ is the required
workflow reliability. Task $i\in\mathcal{N}_k$ has workload $w_i$ measured
in Giga-cycles, and
$d_{ij}$ denotes the data transferred along edge
$(i,j)\in\mathcal{E}_k$. Workflow $k$ arrives at time $a_k$.

The platform contains edge servers $\mathcal{M}^{\mathrm{e}}$ and cloud
servers $\mathcal{M}^{\mathrm{c}}$. Their union is
\begin{equation}
    \mathcal{M}=\mathcal{M}^{\mathrm{e}}\cup\mathcal{M}^{\mathrm{c}}.
    \label{eq:server_set}
\end{equation}
Each server executes one task at a time without preemption. A task becomes
ready after all its predecessors finish. Among ready tasks, a fixed
smallest-workload-first rule determines the next task, while the scheduling
policy selects its server and operating frequency.

\subsubsection{Communication Delay and Execution Time}
Let $m_i\in\mathcal{M}$ be the server selected for task $i$, and let $B_{mn}$
denote the effective bandwidth from server $m$ to server $n$. For edge
$(i,j)\in\mathcal{E}_k$, the communication delay is
\begin{equation}
    c_{ij}(m_i,m_j)=
    \begin{cases}
        0, & m_i=m_j,\\[2pt]
        \dfrac{d_{ij}}{B_{m_i m_j}}, & m_i\neq m_j.
    \end{cases}
    \label{eq:communication_time}
\end{equation}
Server $m$ has maximum processing frequency $f_m^{\max}$ in GHz. For task
$i$, the DVFS ratio $\rho_i\in[\rho_{m_i}^{\min},1]$ determines its actual
frequency as $f_i=\rho_i f_{m_i}^{\max}$, where $\rho_m^{\min}$ is the
minimum DVFS ratio supported by server $m$. The execution time of task $i$ is
\begin{equation}
    t_i=\frac{w_i}{f_i}
       =\frac{w_i}{\rho_i f_{m_i}^{\max}}.
    \label{eq:execution_time}
\end{equation}
Initial data upload and final-result delivery are not included in the
scheduling boundary.

\subsubsection{Online Queue-Based Time Model}
Let $\operatorname{pred}(i)$ denote the predecessor set of task $i$. Its ready
time on candidate server $m$ is
\begin{equation}
    RT_i(m)=
    \begin{cases}
        a_k, & \operatorname{pred}(i)=\varnothing,\\[3pt]
        \displaystyle\max_{j\in\operatorname{pred}(i)}
        \{AFT_j+c_{ji}(m_j,m)\}, & \text{otherwise}.
    \end{cases}
    \label{eq:ready_time}
\end{equation}
Let $T_m$ be the time at which server $m$ becomes available, including work
already queued by other workflows. The start and finish times of task $i$ are
\begin{equation}
    AST_i=\max\{RT_i(m_i),T_{m_i}\},
    \qquad
    AFT_i=AST_i+t_i.
    \label{eq:task_times}
\end{equation}
After assignment, $T_{m_i}$ is updated to $AFT_i$. This recurrence represents
both precedence constraints and non-overlapping execution on each server.

\subsection{Energy Model and DVFS}
\label{sec:energy_model}

Let $P_m^0$ and $P_m^d$ denote the base active power and
maximum-frequency dynamic power of server $m$, respectively. At DVFS ratio
$\rho$, its active power is
\begin{equation}
    P_m(\rho)=P_m^{0}+P_m^{d}\rho^3.
    \label{eq:power_model}
\end{equation}
The computing energy consumed by task $i$ is
\begin{equation}
    E_i=P_{m_i}(\rho_i)t_i.
    \label{eq:task_energy}
\end{equation}
The total computing energy of workflow $G_k$ is
\begin{equation}
    E_k=\sum_{i\in\mathcal{N}_k}E_i.
    \label{eq:total_energy}
\end{equation}
Communication energy and schedule-independent idle energy are excluded.

\subsection{Reliability Model}
\label{sec:reliability_model}

Let $\lambda_m^0$ be the transient-failure rate of server $m$ at maximum
frequency, and let $\beta_m$ describe its sensitivity to frequency reduction.
Its failure rate at DVFS ratio $\rho$ is
\begin{equation}
    \lambda_m(\rho)=\lambda_m^0
    10^{\beta_m(1-\rho)/(1-\rho_m^{\min})}.
    \label{eq:frequency_failure_rate}
\end{equation}
Assuming independent exponential failures, the reliability of task $i$ and
workflow $G_k$ is
\begin{equation}
    R_i=\exp[-\lambda_{m_i}(\rho_i)t_i],
    \qquad
    R_k=\prod_{i\in\mathcal{N}_k}R_i.
    \label{eq:workflow_reliability}
\end{equation}
The reliability requirement is satisfied when $R_k\geq\Theta_k$.

\subsection{Deadline Model}
\label{sec:deadline_model}

The response time of workflow $G_k$ is measured from its arrival to the
completion of its exit task:
\begin{equation}
    C_k=\max_{i\in\mathcal{N}_k}AFT_i-a_k.
    \label{eq:workflow_response_time}
\end{equation}
The deadline requirement is satisfied when $C_k\leq D_k$. The procedure used
to generate $D_k$ for the experiments is described in
Section~\ref{sec:evaluation}.

\subsection{Optimization Problem}
\label{sec:optimization_problem}

For each arriving workflow, the scheduler selects the server--DVFS pair
$(m_i,\rho_i)$ of every task. The problem is formulated as
\begin{subequations}
\label{eq:optimization}
\begin{align}
    \min_{\{m_i,\rho_i\}}
    \quad & E_k
    \label{eq:optimization_objective}\\
    \text{s.t.}\quad
    & C_k\leq D_k,
    \label{eq:constraint_deadline}\\
    & R_k\geq\Theta_k,
    \label{eq:constraint_reliability}\\
    & m_i\in\mathcal{M},\quad
      \rho_i\in[\rho_{m_i}^{\min},1],
    && \forall i\in\mathcal{N}_k.
    \label{eq:constraint_actions}
\end{align}
\end{subequations}
The timing relations in~\eqref{eq:ready_time}--\eqref{eq:task_times} enforce
task precedence and server exclusivity. Since $m_i$ is discrete and $\rho_i$ is
continuous, Problem~\eqref{eq:optimization} has a mixed action space. The next
section develops a hierarchical reinforcement learning method for making these
decisions online.

\section{Proposed Hierarchical Constrained Reinforcement Learning Method}
\label{sec:method}

\subsection{Hierarchical Decision Process}

\subsection{Upper-Level Three-Head Dueling Double DQN}

\subsection{Lower-Level Constrained Soft Actor--Critic}

\subsection{Constraint Credit Assignment and Lagrangian Updates}

\subsection{Constraint-First Checkpoint Selection}

\section{Performance Evaluation}
\label{sec:evaluation}

\subsection{Experimental Setup}

\subsection{Baselines and Implementation Details}

\subsection{Overall Performance}

\subsection{Deadline and Reliability Analysis}

\subsection{Ablation Study}

\section{Conclusion}
\label{sec:conclusion}

\bibliographystyle{IEEEtran}
\bibliography{references}

\end{document}
